\section{Hash Security of the Root Tag}
\label{sec:cmtds3}

The tag-level guarantees of this section rest on one construction, isolated
below as the Root tag Candidate Sequence lemma
(\Cref{lem:root-candidate-seq}): a chronological candidate sequence of
root tag inputs, whose collision or preimage probability the adaptive
lemmas of \Cref{app:lemmas} bound, controlling the chance that distinct
encryption inputs land on the same $\tauroot$-bit root tag projection.
We instantiate it in the random-oracle model for collision and preimage
resistance of the wrapper $\HTW$ and for the \cite{BH22} CMTDs-4 case;
\Cref{app:flat-sponge-lift} transfers each instance to the
public-permutation setting and \Cref{sec:headline} records the resulting
concrete bounds.

First, the random-oracle CMTDs-4 bound (\Cref{thm:cmtds4-ro}) is the
pure $s$-way root tag collision case; the headline CMTDs-4 statement
(\Cref{cor:cmtds4}) is its flat-sponge lift, and
\Cref{app:committing-transfer} forwards the remaining committing notions
to the same root-collision bound. Second, the random-oracle preimage
bound (\Cref{thm:pre-ro}) runs the same argument with the $s$-way
collision replaced by a preimage to a fixed target tag, underlying the
headline preimage statement (\Cref{thm:pre}).
Third, the random-oracle collision bound (\Cref{thm:coll-ro}) keeps
the root tag collision argument intact but, because colliding inputs may
use different ciphertext bodies, first rules out the
only other way distinct bodies can reach a common root transcript---a
collision among the absorbed leaf tags---through the leaf tag collision
bound (\Cref{lem:leaf-collision}); its
public-permutation form is the headline \Cref{thm:coll}.

The committing and discoverability arguments below rest on one structural
fact, which we record first: distinct contexts reach distinct final root
transcripts before any leaf tag is aggregated, so agreement among leaf tags
cannot merge them.

\begin{proposition}[Opening Triple Separation]
\label{cor:triple-sep}
If two \TW{} root computations have distinct $(K, N, A)$ triples, then
their bridged final root transcript inputs under $\flatmap(\cdot)$ are
distinct, regardless of whether each computation is an encryption or
a decryption check on a common ciphertext.
\end{proposition}

\begin{proof}
Suppose for contradiction that $\flatmap(R_i) = \flatmap(R_j)$. Since the
bridge $\flatmap(\cdot)$ is injective on well-formed transcript prefixes
(\Cref{lem:bridge-decode}), the native flattened inputs of the two
final-root prefixes are equal. Then \Cref{lem:transcript-inj}
recovers the full duplex-call sequence and reads off the pair
$(\sigma,\sigph)$ for each call. The seven $4$-bit phase suffixes are
pairwise distinct, so the \siginit, $\sigma_{\mathsf{ad}}^\pm$,
$\sigma_{\mathsf{msg}}^\pm$, and (when present) $\sigma_{\mathsf{agg}}^\pm$
calls in the recovered sequence are unambiguously identified. The
aggregation phase is present exactly when $n \geq 1$, and
$n = \leaves(C)$ is determined by the common ciphertext, so the two
recovered sequences either both carry an aggregation phase or both omit
it; either way the recovery of $(K, N, A)$ below uses only the
\siginit and AD calls, which precede any message or aggregation block.

Equality of the recovered sequences forces equality of the \siginit
call, whose $\sigma$ is $\prfxroot \concat K \concat N$ with
fixed-width $K$ and $N$, hence $K_i = K_j$ and $N_i = N_j$. It also
forces equality of the recovered AD subsequences. By the AD bullet of
\Cref{lem:transcript-inj} the associated-data phase contributes a
$\sigadlast$-tagged block exactly when the associated data is nonempty,
so the recovered sequences agree on the presence of the AD phase. If
exactly one of $A_i, A_j$ were empty the recovered sequences would
differ in the presence of a $\sigadlast$ call, contradicting their
equality; hence either both are empty, giving $A_i = A_j = \varepsilon$,
or both are nonempty, in which case the phase-recovery sub-claim of
\Cref{lem:transcript-inj} inverts the $\rho$-byte-block partition map on
the AD byte string and equal AD $\sigma$-block sequences imply
$A_i = A_j$. In all cases $(K_i, N_i, A_i) = (K_j, N_j, A_j)$,
contradicting the hypothesis of distinct triples.
\end{proof}

\begin{lemma}[Root tag Candidate Sequence]
\label{lem:root-candidate-seq}
Consider a random-oracle game of \Cref{sec:ro-model} in which the
adversary $\mathcal{B}$ first makes direct $\ROflat$ queries and then
outputs values on which the challenger runs $m \geq 1$ deterministic
encryption or decryption checks, the $i$-th reaching a final root transcript $R_i$
whose root tag is read by a single $\ROflat$ lookup at $\flatmap(R_i)$.
Every execution then admits a chronological sequence of root tag inputs
such that
\begin{enumerate}
\item each distinct direct query whose input $Y$ decodes under
  $\KeyedTWOut$ to a root tag output-point class---$\Rmsg^+$ (the final
  root message block, which carries the root tag only on the no-leaf path)
  or $\Rtag$ (the aggregated root tag)---is appended as a candidate before
  its $\ROflat$ lookup, and contributes at most one candidate; multi-chunk
  computations also reach $\Rmsg^+$ with zero requested output before
  aggregation, so this predicate over-approximates the root tag candidate
  set, which only loosens the count;
\item if the challenger reaches its checks, each $\flatmap(R_i)$ is in
  the sequence, appended---if not already present---immediately before
  the $i$-th root tag answer is sampled;
\item no candidate's $\ROflat$ value is read before its append step, so
  the sequence satisfies the predictability and unrevealed-value premises
  of \Cref{lem:adaptive-candidate};
  and
\item its length is at most $\qRO(\mathcal{B}) + m$.
\end{enumerate}
\end{lemma}

\begin{proof}
The committing, discoverability, and hash (collision, second-preimage,
and preimage) games are keyless: all keys are
adversary-supplied, so a direct query to a keyed transcript is one of
$\mathcal{B}$'s visible $\ROflat$ queries counted in $\qRO(\mathcal{B})$,
and there is no $\badkey$ event. No construction oracle runs during the
direct-query phase, so each direct-query candidate is first read no
earlier than its append step. The append predicate depends on $Y$ alone
through the exact decoder of \Cref{lem:bridge-decode}, which ignores the
requested output length, and by Output-Point Separation
(\Cref{cor:output-sep}) each input decodes to at most one output-point
class; hence each distinct direct query contributes at most one
candidate, giving~(1).

If the challenger reaches its $m$ checks, append $\flatmap(R_i)$
immediately before the $i$-th root tag answer is sampled, unless it is
already present. The intermediate $\ROflat$ lookups of check $i$ lie at
non-root tag output points, so by \Cref{cor:output-sep} they do not read
$\flatmap(R_i)$. If $\flatmap(R_i)$ was touched earlier---by a direct
query or by an earlier check $j < i$---then \Cref{lem:bridge-decode} and
\Cref{cor:output-sep} make that earlier access a final-root lookup at the
same address, which was therefore already appended, so the input is in
the sequence and is not appended again. Otherwise its value is still
unread when appended. Either way no candidate's value is read before its
append step, giving~(2) and~(3). At most $\qRO(\mathcal{B})$ direct
candidates and $m$ challenger candidates are appended, giving~(4).
\end{proof}

\begin{theorem}[Random-Oracle CMTDs-4]
\label{thm:cmtds4-ro}
For every $s \geq 2$ and every $s$-way CMTDs-4
adversary $\mathcal{B}$ in the \TW{} random-oracle model
(\Cref{sec:ro-model}),
\[
  \Advcmtds{4}_\RO(\mathcal{B})
  \;\leq\; \RootColl_{\tauroot}(\qRO(\mathcal{B}), s)
  = \frac{\binom{\qRO(\mathcal{B}) + s}{s}}{2^{\tauroot (s - 1)}}.
\]
\end{theorem}

\begin{proof}
Run the random-oracle CMTDs-4 game (\Cref{sec:ro-model}). After
rejecting on any failed syntax, message-length, or pairwise-full-input
check, its challenger performs $m = s$ deterministic decryption checks
$\Dec_{K_i}^\RO(N_i,A_i,C \concat T)$ on the one common body
$C \concat T$, reaching final root transcripts $R_1,\dots,R_s$. The
Root tag Candidate Sequence lemma (\Cref{lem:root-candidate-seq}) yields
a sequence of length at most $\qRO(\mathcal{B}) + s$ meeting the premises
of \Cref{lem:adaptive-candidate} and containing each
$\flatmap(R_i)$.

On a winning execution the pairwise-full-input check passed, so the
tuples $(K_i,N_i,A_i,M_i)$ are pairwise distinct. Because all $s$ checks
decrypt the common body $C \concat T$, determinism of $\Dec$ makes the
$(K_i,N_i,A_i)$ triples pairwise distinct as well: equal triples would
return equal messages, hence equal tuples. By Opening Triple Separation
(\Cref{cor:triple-sep}), the bridged final-root inputs
$\flatmap(R_1),\dots,\flatmap(R_s)$ are therefore pairwise distinct; leaf
tag collisions cannot merge them. Since every decryption check accepts,
the $\tauroot$-bit projections at these $s$ distinct candidates all equal
the supplied tag $T$, so the win event is contained in the $s$-way
candidate-collision event. \Cref{lem:adaptive-candidate} with
$L=\qRO(\mathcal{B})+s$ and $\tau=\tauroot$ gives
\[
  \Advcmtds{4}_\RO(\mathcal{B})
  \leq
  \binom{\qRO(\mathcal{B}) + s}{s} \cdot 2^{-\tauroot(s-1)}
  =
  \RootColl_{\tauroot}(\qRO(\mathcal{B}), s).
\]
\end{proof}

\begin{theorem}[Random-Oracle Collision Resistance of $\HTW$]
\label{thm:coll-ro}
For every $s \geq 2$ and every $s$-way collision
adversary $\mathcal{B}$ in the \TW{} random-oracle model,
\[
  \AdvColl_{\RO}(\mathcal{B})
  \;\leq\;
  \LeafColl_{\tauleaf}(\qRO(\mathcal{B})+\nleaf(\mathcal{B}))
  + \RootColl_{\tauroot}(\qRO(\mathcal{B}), s).
\]
\end{theorem}

\begin{proof}
Run the random-oracle collision game. Let $\badleaf^\star$ be the event
that two distinct construction-generated final leaf transcripts receive
the same $\tauleaf$-bit $\RF$ projection. By the adversary-supplied-key
bound of \Cref{lem:leaf-collision}, $\Pr[\badleaf^\star] \leq
\LeafColl_{\tauleaf}(\qRO(\mathcal{B})+\nleaf(\mathcal{B}))$, the stated
leaf-collision term, so it remains to bound the probability of the win
event under $\neg\badleaf^\star$.

After rejecting unless the $s$ output inputs are pairwise distinct and in
the \TW{} domain, the challenger performs $m = s$ deterministic
encryptions $\Enc_{K_i}^{\RO}(N_i,A_i,M_i)$, reaching
final root transcripts $R_1,\dots,R_s$. The Root tag Candidate Sequence
lemma (\Cref{lem:root-candidate-seq}) yields a sequence of length at most
$\qRO(\mathcal{B})+s$ meeting the premises of
\Cref{lem:adaptive-candidate} and containing each
$\flatmap(R_i)$.

On a winning execution with $\neg\badleaf^\star$, the inputs
$(K_i,N_i,A_i,M_i)$ are pairwise distinct and all $s$ encryptions yield a
common tag $T$. We first show that the bridged final
root inputs $\flatmap(R_1),\dots,\flatmap(R_s)$ are pairwise distinct.
Suppose instead that checks $i \neq j$ reach the same final root input.
The bridge and transcript encodings are injective, so the root
transcripts have the same key, nonce, associated data, root ciphertext
chunk, root tag output point, and, when leaves are present, the same
ordered leaf tag aggregation string. The trailing encoded leaf count
gives the same number of leaves. For each leaf position, equality of the
aggregated leaf tag values together with $\neg\badleaf^\star$ implies
that the two final leaf transcripts are the same; hence the
corresponding leaf chunks are the same. Thus the two encryptions share
$(K,N,A,M)$, so $(K_i,N_i,A_i,M_i)=(K_j,N_j,A_j,M_j)$, contradicting the
pairwise distinctness of the colliding inputs. The final root inputs are
therefore pairwise distinct.

Since all $s$ encryptions yield the common tag, the $\tauroot$-bit
projections at these $s$ distinct final-root candidates are all equal. The win event under $\neg\badleaf^\star$ is contained in the
$s$-way root candidate-collision event, which
\Cref{lem:adaptive-candidate} bounds by
$\RootColl_{\tauroot}(\qRO(\mathcal{B}),s)$. Adding the
$\badleaf^\star$ probability gives the theorem.
\end{proof}

\begin{theorem}[Random-Oracle Preimage Resistance of $\HTW$]
\label{thm:pre-ro}
For every query-independent target tag $T^* \in \bits^{\tauroot}$ and
every preimage adversary $\mathcal{B}$ in the \TW{}
random-oracle model (\Cref{sec:ro-model}),
\[
  \AdvPre_\RO(\mathcal{B})
  \;\leq\; \RootPre_{\tauroot}(\qRO(\mathcal{B}))
  = \frac{\qRO(\mathcal{B}) + 1}{2^{\tauroot}}.
\]
\end{theorem}

\begin{proof}
Run the random-oracle preimage game (\Cref{sec:ro-model}) for the
query-independent target $T^*$, which is independent of $\ROflat$. After
rejecting unless the output input $X = (K,N,A,M)$ is in the \TW{} domain,
the challenger performs $m = 1$ deterministic encryption
$\Enc_{K}^\RO(N,A,M)$, reaching a final root transcript $R$. The Root tag
Candidate Sequence lemma
(\Cref{lem:root-candidate-seq}) yields a sequence of length at most
$\qRO(\mathcal{B}) + 1$ meeting the premises of
\Cref{lem:adaptive-candidate} and containing $\flatmap(R)$; with
a single input no pairwise distinctness is needed.

On a winning execution $\HTW(X) = T^*$, that is, the $\tauroot$-bit
projection $T'$ at the final root transcript $R$ equals $T^*$; in the
random-oracle model
$T' = \bigl\lfloor \ROflat(\flatmap(R)) \bigr\rfloor_{\tauroot}$. Hence
the win event is contained in the event that some candidate's
$\tauroot$-bit projection equals the fixed target $T^*$. Applying
\Cref{lem:adaptive-candidate} with $L = \qRO(\mathcal{B}) + 1$,
$\tau = \tauroot$, and $t = T^*$ gives
\[
  \AdvPre_\RO(\mathcal{B})
  \leq
  (\qRO(\mathcal{B}) + 1) \cdot 2^{-\tauroot}
  =
  \RootPre_{\tauroot}(\qRO(\mathcal{B})).
\]
\end{proof}

\begin{theorem}[Random-Oracle aSec]
\label{thm:asec-ro}
For every aSec adversary $\mathcal{B}$ in the \TW{} random-oracle model
(\Cref{sec:ro-model}),
\[
  \AdvaSec_\RO(\mathcal{B})
  \;\leq\;
  \RootPre_{\tauroot}(\qRO(\mathcal{B}))
  + \LeafColl_{\tauleaf}(\qRO(\mathcal{B})+\nleaf(\mathcal{B})).
\]
\end{theorem}

\begin{proof}
Run the random-oracle aSec game. The challenger draws the challenge input
$X$, encrypts it to obtain $T = \HTW(X)$ at a final root transcript $R_X$,
and gives $X$ to $\mathcal{B}$, which outputs an input $X' \neq X$; the
win event is $\HTW(X') = T$. By the adversary-supplied-key bound of
\Cref{lem:leaf-collision} the probability of a collision
$\badleaf^\star$ among the leaf transcripts of the encryptions of $X$ and
$X'$ is at most the stated leaf term, so it remains to bound the win
under $\neg\badleaf^\star$.

Here $T = \lfloor \RF(R_X) \rfloor_{\tauroot}$ is the root tag of the
challenge transcript $R_X$. Encrypting $X'$ reaches a final root
transcript $R'$ with $\HTW(X') = \lfloor \RF(R') \rfloor_{\tauroot}$. If
$R' = R_X$, then by transcript injectivity (\Cref{lem:transcript-inj})
$X'$ shares the context, leaf count, and aggregated leaf tags of $X$, and
under $\neg\badleaf^\star$ equal leaf tags force equal leaf transcripts
by Leaf-Transcript Recovery (\Cref{lem:leaf-recovery}),
hence equal leaf bodies; so $X' = X$ by determinism, contradicting
$X' \neq X$. Thus $R' \neq R_X$, and the win event $\HTW(X') = T$ is the
event that $R'$ collides with $R_X$ on its $\tauroot$-bit projection. The
Root tag Candidate Sequence lemma (\Cref{lem:root-candidate-seq}), applied
to $\mathcal{B}$'s direct queries and the single post-output check
encrypting $X'$, yields a chronological sequence of at most
$\qRO(\mathcal{B})+1$ candidates containing $\flatmap(R')$, each appended
with unrevealed value; prepending the challenge candidate $\flatmap(R_X)$
as the designated position---whose value is the exposed challenge tag
$T$---the second-preimage corollary of \Cref{lem:adaptive-candidate} bounds
the probability that any candidate other than $\flatmap(R_X)$ collides
with it by $\RootPre_{\tauroot}(\qRO(\mathcal{B}))$. Adding the
$\badleaf^\star$ probability gives the theorem.
\end{proof}
